\section{Data and Methods}\label{sec:data-methods}

\subsection{Census scope and inclusion criteria}\label{dm:census}
We census open-weight, decoder-only language models published on the Hugging Face Hub
between January 2023 and the census snapshot in October 2026. We keep checkpoints released
by the original publisher organization (\kvcOrgCount{} organizations) and exclude
fine-tunes, merges, and quantized re-uploads from third-party accounts. Multimodal models
enter through their text backbone: we parse the language-model sub-configuration and ignore
vision and audio towers. The census holds \kvcCensusRows{} checkpoint rows ---
\kvcYearRowsTwentyThree{} released in 2023, \kvcYearRowsTwentyFour{} in 2024,
\kvcYearRowsTwentyFive{} in 2025, and \kvcYearRowsTwentySix{} in 2026 to date. Every
\texttt{config.json} is snapshotted and all Hub API responses are cached on disk, so the
pipeline re-runs offline and deterministically.

Selection effects are counted, not hidden. Gated repositories whose \texttt{config.json}
is not publicly downloadable are skipped; \kvcGatedSkipped{} repositories were skipped this
way, which biases the census away from gated families and is reported as a selection-bias
limitation. The release date is the Hub \texttt{createdAt} timestamp, a proxy that can
predate or lag the effective public release. Parsing failed for \kvcKvErrRows{}
configurations; these are listed in \texttt{data/kv\_config\_errors.txt} and excluded from
the census table. \kvcNoParams{} included rows carry no parameter count (neither
safetensors metadata nor a sufficient configuration); both gaps are carried as
limitations. Where a count exists, it comes from
Hub safetensors metadata, otherwise it is computed from the configuration; the per-row
source is recorded in the \texttt{params\_source} column.

\subsection{KV-cache accounting conventions}\label{dm:kv}
All cache sizes are analytic per-sequence minima computed by \texttt{analysis/kv.py}, a
library of pure functions over the snapshotted configuration dictionaries; it never loads
weights. A full-attention layer caches one key and one value tensor of
$n_{\mathrm{kv}} \times d_h$ elements per token, i.e.\ $2\,n_{\mathrm{kv}}\,d_h\,B$ bytes
per token, where $B$ is the bytes per element (2 for BF16). Grouped-query and multi-query
attention are the same formula with reduced $n_{\mathrm{kv}}$
\cite{shazeer2019mqa,ainslie2023gqa}. A sliding-window layer with window $W$ holds
$\min(T,W)$ tokens at context $T$ \cite{jiang2023mistral}. A multi-head latent attention
layer caches only the compressed latent, $(r_{\mathrm{kv}}+d_{\mathrm{rope}})\,B$ bytes
per token, where $r_{\mathrm{kv}}$ is the latent rank (\texttt{kv\_lora\_rank}) and
$d_{\mathrm{rope}}$ the decoupled RoPE dimension (\texttt{qk\_rope\_head\_dim})
\cite{deepseekai2024v2}; values are decompressed on the fly, so we count the theoretical
latent and measure the as-implemented cache in the validation experiment. A recurrent
layer --- Mamba-style state-space layers \cite{gu2023mamba}, linear attention
\cite{katharopoulos2020linear}, gated delta-rule, RWKV --- holds a constant state whose
size is computed from the configuration (SSM and convolution state, or a matrix state) and
does not depend on $T$. Falcon-H1 layers carry both an attention and a recurrent cache,
and each is counted. Caches shared across layers are counted once \cite{sun2024yoco}.

The headline outcome is the asymptotic marginal cache cost per token,
$\lim_{T\to\infty}\mathrm{d}\,\mathrm{bytes}(T)/\mathrm{d}T$: only layers whose cache grows
without bound (full and latent layers) contribute, since windowed layers saturate at $W$
tokens and recurrent state is constant. Labels are honest by construction: a model is
labelled \texttt{mha}, \texttt{gqa}, or \texttt{mqa} only if it is pure full attention and
the computed per-token bytes equal $2\,L\,n_{\mathrm{kv}}\,d_h\,B$ exactly; any windowed,
recurrent, latent, or shared component forces a descriptive label such as
\texttt{gqa+swa}, \texttt{hybrid-mamba}, or \texttt{mla}. Family-specific parsing rules
follow the upstream \texttt{transformers} configuration semantics --- explicit head
dimensions, per-family window flags, hybrid layer-pattern strings --- cross-checked against
published configurations \cite{grattafiori2024llama3,yang2025qwen3}. The library carries
107 unit tests against hand-computed values and 61 real configuration snapshots; an
unknown architecture raises an error rather than defaulting to a guess.

\subsection{Pre-registered analysis}\label{dm:plan}
The analysis plan (\texttt{analysis/PLAN.md}) was fixed before any trend statistic was
computed, and every deviation is recorded there as a dated note. Trend and regression
statistics use a deduplicated set: checkpoints with identical attention geometry share a
\texttt{dedup\_group}, and the representative is the earliest-released member, ties broken
by repository id. This leaves \kvcDedupRows{} rows in \kvcGroups{} groups. RQ1 reports
per-year medians and interquartile ranges of the two primary outcomes and Spearman rank
correlations against release date, with pre-registered support thresholds. RQ2 regresses
$\log_{10}$ bytes/token on $\log_{10}$ parameters by OLS, with 95\% confidence intervals
from 1000 bootstrap resamples at seed 0, and fits a mechanism regression adding GQA ratio,
head dimension, depth, window share, recurrent share, a latent indicator, and a
shared-cache indicator, reported as standardized coefficients with variance inflation
factors. The MoE-versus-dense comparison uses a two-sided Mann-Whitney $U$ test with a
rank-biserial effect size. Capacity computations use all rows: each checkpoint is its own
capacity point.

\subsection{Capacity model}\label{dm:capacity}
For every checkpoint we solve a one-dimensional memory budget: usable KV pool $=$ budget
$-$ weight bytes $-$ runtime overhead. Budgets are 24, 80, and 128 GiB ($2^{30}$ bytes),
matching common accelerator and unified-memory tiers. Weights occupy 2 bytes per parameter
(BF16), 1 byte (FP8), or 4.5 bits per parameter under 4-bit quantization, the effective
figure including 16 group-scale bits per 32-element group \cite{frantar2022gptq}. The KV
cache is priced at 2 or 1 byte per element (BF16 or FP8) \cite{liu2024kivi}; we do not
model allocator fragmentation, which serving systems manage separately
\cite{kwon2023paged}. Runtime overhead is a free parameter: we report the full grid over
0, 2, 6, and 12 GiB and use 6 GiB in the headline tables. Because $\mathrm{bytes}(T)$ is
monotone piecewise linear with breakpoints at the window sizes, the maximum token count
under the pool and the KV-parity context $T^{*}$ --- the length at which one sequence's
cache equals the BF16 weight bytes --- are solved exactly on each linear segment. The plan
pre-registered a log-spaced grid with interpolation; the exact solve is a strict
refinement of the same model, recorded as a dated deviation.

\subsection{Validation against allocated caches}\label{dm:validation}
Analytic minima are meaningless if the reference stack allocates something else. For RQ4
we instantiate reduced models directly from the snapshotted configurations with
\texttt{from\_config} (random initialization, no weight downloads, CPU only). Depth is cut
to the smallest value that preserves one full period of the layer-type pattern, the
vocabulary is shrunk, and per-layer attention geometry --- KV heads, head dimensions,
latent ranks, window sizes --- is kept exact. We prefill $T$ tokens, with $T>W$ for
windowed models, and sum the bytes of every cache tensor the returned cache object holds:
keys and values, convolution and SSM states. The analytic side is evaluated on the same
reduced configuration in float32. The relative error is
$\lvert\mathrm{measured}-\mathrm{analytic}\rvert/\mathrm{analytic}$; every row above 1\%
must carry an explanation, and the explanation itself is a reported finding. The
experiment covers \kvcValRows{} models spanning \kvcValClasses{} attention classes, with
\kvcValWithin{} rows within 1\%.
